%% cover-letter.tex  —  The American Mathematical Monthly
%% Submission portal: Taylor & Francis ScholarOne
%% Author: Ali Elouafiq, SQLI Digital Laboratory, Levallois-Perret, France.
%% =============================================================================
\documentclass[12pt]{article}
\usepackage[margin=1in]{geometry}
\usepackage[T1]{fontenc}
\usepackage{hyperref}
\pagestyle{empty}

\begin{document}

\noindent
Ali Elouafiq\\
Principal Research Engineer\\
SQLI Digital Laboratory\\
2--10 rue Thierry Le Luron\\
92300 Levallois-Perret, France\\
\href{mailto:ali@sqli.com}{ali@sqli.com}

\bigskip
\noindent
9 June 2026

\bigskip
\noindent
Editor-in-Chief\\
\textit{The American Mathematical Monthly}\\
Mathematical Association of America

\bigskip
\noindent
Dear Editor,

\medskip
\noindent
I submit for your consideration the manuscript:

\begin{center}
  \textbf{The continuous-coefficient Jensen equation:}\\
  \textbf{A note on vestigial regularity hypotheses}
\end{center}

\noindent
for publication in \textit{The American Mathematical Monthly} as an
\textbf{Article}.

\bigskip
\noindent\textbf{What the paper does.}
The paper shows that the continuous-coefficient Jensen equation
$G(pu_1+(1-p)u_2)=pG(u_1)+(1-p)G(u_2)$ for all $u_1,u_2\in[0,M]$ and
all $p\in[0,1]$ forces $G$ to be affine with \emph{no regularity
hypothesis}: a one-line endpoint substitution ($u_1{:=}M$, $u_2{:=}0$,
$p{:=}v/M$) reads $G(v)$ directly off the straight line through the
endpoints. The paper accompanies this one-line observation with
(i)~a dictionary comparing the result with the well-known case of the
discrete-coefficient Jensen equation (where Hamel's pathology is real
and regularity \emph{is} needed), (ii)~a structural explanation of
why the two cases differ, and (iii)~three application areas\,---\,%
expected-utility theory, the axiomatic characterization of Shannon
entropy, and surrogate-loss calibration\,---\,where the same
vestigial hypotheses recur.

\bigskip
\noindent\textbf{Why the Monthly.}
The result is elementary enough for a second-year analysis student to verify
in an afternoon, yet it is consistently missed in practice: the functional
equations literature treats the two forms of Jensen's equation in adjacent
sections without flagging the asymmetry, and modern applied work (utility
theory, calibration, entropy) carries defensive regularity assumptions that
are unnecessary precisely in the continuous-coefficient case.  The Monthly
is the natural home for an expository note that makes this asymmetry
explicit, provides the dictionary, and invites readers to audit their own
work.

\bigskip
\noindent\textbf{Novelty and relationship to prior art.}
The substitution technique that closes Theorem~1 is folklore:
Acz\'{e}l's 1966 monograph (\S2.1.4) contains a dyadic precursor proved
under a continuity hypothesis, and the genre is what Kuczma's 2009
introduction describes as the part of Jensen-equation theory that does
not require regularity assumptions on the function. The precise
continuous-coefficient regularity-free statement of Theorem~1 does
not, however, appear in this exact form in the standard references the
author consulted. The contribution of this note is therefore not
Theorem~1 itself but its \emph{explicit packaging}: the dictionary
contrasting the two Jensen forms, the structural explanation of the
Hamel-pathology mechanism, and the identification of a recurrence
pattern in derivations that produce the continuous-coefficient form as
a saturated identity. A closing note records that the observation
reached the author through a formalization project on a separate
topic\,---\,a route that may interest readers curious about how
mechanized proof can pressure-test classical heuristics.

\bigskip
\noindent\textbf{Submission details.}
\begin{itemize}\setlength\itemsep{2pt}
  \item \textbf{Manuscript type:} Article (expository).
  \item \textbf{Page count:} 9 pages anonymous body, 10 pages with title page.
  \item \textbf{Word count (approximate):} 3{,}400 words (TeX-counted body).
  \item \textbf{MSC 2020 primary:} 39B22 (Functional equations for real functions).
  \item \textbf{MSC 2020 secondary:} 26B25 (Convexity), 91B16 (Utility theory),
        94A17 (Measures of information; entropy).
  \item \textbf{Double-blind:} Yes. The anonymous body
        (\texttt{manuscript-anon.pdf}) contains no author-identifying
        information; the file with author details
        (\texttt{manuscript.pdf}) carries the title page on page~1.
  \item \textbf{Preprint status:} The manuscript has not been deposited on
        arXiv. It is not under consideration elsewhere.
  \item \textbf{Competing interests:} None.
  \item \textbf{Funding:} This research received no specific grant from any
        funding agency in the public, commercial, or not-for-profit
        sectors.
  \item \textbf{Use of generative AI:} Disclosed in the manuscript's
        Acknowledgments section. Anthropic's Claude Opus~4.8 was used
        for adversarial review of drafts and \LaTeX{} formatting; all
        mathematical content was conceived, verified, and refined by
        the author.
\end{itemize}

\bigskip
\noindent
I look forward to the referees' comments.

\bigskip
\noindent
Sincerely,

\vspace{28pt}
\noindent
Ali Elouafiq
\end{document}
